\documentclass[12pt]{article}
\usepackage{graphicx}
\usepackage{amsmath}
\usepackage{url}
\usepackage{booktabs}
\usepackage[margin=1in]{geometry}
\usepackage{natbib}

\title{The Revenue Premium at the Border: State Cannabis Policy Gaps and\\
Dispensary Sales in the United States}
\author{Mahyar Rezazadeh \\ Babak Sahragard}
\date{August 2026}

\begin{document}
\maketitle

\begin{abstract}
\noindent We ask three questions about cannabis retail markets in U.S.\ states located near a state where
recreational sales remain illegal. First, does a border premium exist: is revenue per dispensary higher in
counties near a nonlegal-state border than in interior counties of the same state, using all border counties
as treatment and all interior counties as control? Using a county-year panel for Colorado and Washington, we
find it does, both with a 70-mile binary border definition (county-clustered $p=0.033$) and a continuous
distance measure (county-clustered $p=0.023$), while total revenue per capita and a non-cannabis placebo
sector show no comparable pattern. Second, is this evidence of license-cap policy: do states that restrict
entry show the premium on a different margin than states that do not? Testing this directly with Montana, the
one sample state without a statewide numerical cap, we find the reverse pattern there: a significant
per-capita effect and a precise per-dispensary null. Third, how is this demand captured: does the premium
scale with the size of the population on the nonlegal side, not just proximity to it? We show it does
(county-clustered $p=0.036$). Together, these three results establish a border premium, connect it to a
specific and testable state policy mechanism, and show it responds to the size of the market it draws from.
\end{abstract}

\section{Introduction}

Recreational marijuana is legal in a growing number of U.S.\ states but remains prohibited in many others.
By the end of 2025, 24 states had legalized recreational marijuana, although Virginia had not yet
established a legal retail market. Thus, 23 states had operating recreational cannabis markets, while
recreational retail sales remained prohibited in the other 27 states.

This creates sharp policy differences across state borders. Residents of counties in nonlegal states near a
legal market may be able to reach a licensed dispensary with a short drive. Those living farther from
legal-state borders or dispensaries face higher travel costs and have less access to legal recreational
cannabis.

Similar border effects are common when neighboring jurisdictions tax, regulate, or prohibit the same good
differently. Studies of alcohol, cigarettes, gasoline, lotteries, and cannabis show that consumers cross
borders in response to differences in taxes, prices, and legal access. These responses generally weaken with
distance because travel becomes more costly \citep{AsplundFribergWilander2007,Lovenheim2008,BeattyLarsenSommervoll2009,ChandraHeadTappata2014,HansenMillerWeber2020}.
Cross-border sales can also be larger when more potential consumers live near the neighboring market
\citep{KnightSchiff2012,ChandraHeadTappata2014}.

Understanding whether cross-border demand appears in dispensary revenue matters for both state and local
policy. Legal states make decisions about licensing, zoning, and taxation, and these decisions depend partly
on whether dispensaries serve only local residents or also consumers from neighboring nonlegal states.
Nonlegal states may also face a fiscal externality if some consumption and tax revenue shift across the
border. For local governments in border counties, cross-border demand can affect the appropriate number of
dispensaries, traffic, and enforcement needs.

This paper asks three specific questions, organized around three results. \textbf{Result 1} asks whether a
border premium exists at all: using every county in Colorado and Washington within 70 miles of a nonlegal
border as treatment and every other county in those states as control, we test whether revenue per active
dispensary is higher near the border, using both a binary threshold and a continuous distance measure.
\textbf{Result 2} asks whether this reflects license-cap policy specifically: if the number of dispensaries a
state allows cannot adjust to demand, a demand shock should appear as revenue per existing dispensary rather
than as more dispensaries or more revenue per capita. We test this using Montana, the one state in our data
without a statewide numerical cap, as a direct within-sample contrast. \textbf{Result 3} asks how this demand
is captured: whether the premium in Colorado and Washington scales with the size of the nonlegal-side
population within reach, not merely with proximity to it.

The closest study is \citet{HansenMillerWeber2020}, who use cannabis transaction data from Washington and
Oregon's 2015 retail-market opening to show that sales fell sharply at Washington dispensaries near the
Oregon border once Oregon consumers gained legal access at home. This provides direct evidence of
cross-border cannabis shopping. \citet{HaoCowan2020} compare counties directly bordering legal states with
counties farther inside neighboring nonlegal states, but they study marijuana-related arrests rather than
dispensary revenue. Neither paper asks whether a border premium is general across many borders rather than
one policy transition, whether it depends on state licensing policy, or whether it scales with the size of
the population it draws from; these are the three questions this paper takes up in turn.

The remainder of the paper proceeds as follows. Section 2 reviews the relevant literature. Section 3
develops a conceptual framework. Section 4 describes the data. Section 5 presents the empirical strategy.
Sections 6 through 8 report Results 1, 2, and 3 in turn. Section 9 reports robustness checks, Section 10
discusses policy implications and limitations, and Section 11 concludes. A secondary, exploratory analysis
that is not part of the paper's three main results is reported in the appendix.

\section{Literature Review}

A large literature shows that consumers respond to differences in taxes, prices, and regulations across
nearby jurisdictions by shopping across borders. \citet{Lovenheim2008}, for example, finds substantial
cross-border cigarette purchasing in the United States. Similar behavior has been documented for alcohol and
gasoline. \citet{AsplundFribergWilander2007} show that alcohol sales respond to prices in neighboring
countries and that this response becomes weaker with distance from the border. \citet{Stehr2007} finds that
differences in alcohol taxes and Sunday sales restrictions across U.S.\ states also generate cross-border
purchases, while \citet{BanfiFilippiniHunt2005} document a similar pattern for gasoline.

Geography plays an important role in these decisions. \citet{ChandraHeadTappata2014} show that cross-border
shopping falls with distance because consumers must weigh potential savings against travel costs.
\citet{KnightSchiff2012} show that the size of the nearby population also matters. In their study of state
lotteries, cross-border sales are stronger where more consumers live close enough to reach the competing
market at relatively low cost. \citet{Agrawal2015} similarly finds that the effects of state tax differences
are strongest near state borders and weaken with driving time. \citet{BakerJohnsonKueng2021} provide further
evidence that consumers shift purchases across jurisdictions in response to tax differences. Together, these
studies suggest that cross-border demand should depend on both access to the neighboring market and the
number of consumers who can reach it.

These ideas apply naturally to recreational cannabis, where neighboring states can differ not only in taxes
or prices but also in whether recreational sales are legally available. \citet{HansenMillerWeber2020} provide
the most direct evidence of cross-border cannabis shopping. They use detailed transaction data from
Washington and the opening of Oregon's recreational market in 2015. Once Oregon consumers gained access to
stores in their own state, sales at Washington dispensaries near the Oregon border fell sharply, while the
decline was much smaller farther from the border. Their results show that consumers had been crossing state
lines to purchase cannabis. \citet{HaoCowan2020} examine the other side of the policy border. They compare
counties in nonlegal states that directly border Colorado or Washington with counties farther away in those
same states. They find that recreational legalization increased marijuana-possession arrests in border
counties relative to interior counties. Other studies find spillovers in outcomes such as crime and traffic
safety \citep{WuBoatengLang2020,Moyer2026}. Taken together, this evidence shows that the effects of state
cannabis policies can extend beyond the states that adopt them.

The structure of legal cannabis markets is central to Result 2. \citet{HollenbeckUetake2021} show that
restrictions on retail entry in Washington gave existing dispensaries substantial local market power.
\citet{MacePatelSeegert2020} likewise show that competition affects the incidence of cannabis taxes. Local
regulation can create additional constraints because cities and counties may restrict or prohibit cannabis
businesses even when recreational sales are legal at the state level \citep{Yaskewich2022}. These
restrictions matter for our outcome. When entry is limited, an increase in demand does not necessarily lead
to more dispensaries. Existing stores may instead serve the additional consumers and earn higher revenue per
store. This distinction follows the broader market-entry literature, where market size can affect both the
number of firms and the revenue available to each firm \citep{BresnahanReiss1991}.

Cannabis policy may also affect demand for other goods. This is particularly relevant when choosing
comparison sectors for Result 1. Several studies find evidence that cannabis and alcohol can act as
substitutes. \citet{BaggioChongKwon2020} find lower alcohol sales in counties covered by medical marijuana
laws relative to nearby counties without such laws. \citet{AndersonHansenRees2013} also find evidence
consistent with substitution between marijuana and alcohol. \citet{MillerSeo2021} estimate that cannabis
legalization reduced demand for both alcohol and cigarettes in Washington and show that these changes can
affect tax revenue. Other studies find weaker or different relationships. \citet{WenHockenberryCummings2015},
for example, find an increase in binge drinking following medical marijuana legalization, while
\citet{Veligati2020} find no significant overall change in alcohol or cigarette sales.
\citet{AndersonRees2023} conclude, in a comprehensive review, that the relationship depends in part on the
policy and population being studied. For this reason we rely on a placebo sector with no plausible
consumption relationship with cannabis, general retail trade, rather than alcohol.

\section{Conceptual Framework}

We use a simple framework to explain why dispensaries near nonlegal states may earn more revenue. Consider
county $i$ in a state with a legal recreational-marijuana market. Let $d_i$ denote its distance to the
nearest nonlegal state. Residents on the nonlegal side have three main options: they can forgo recreational
cannabis, obtain it through an illegal market, or travel to a legal dispensary across the state border. The
attractiveness of the third option depends partly on how far they must travel and how many potential
consumers live near the legal market.

Let $P_i$ denote the population on the nonlegal side that lives within reasonable reach of dispensaries in
county $i$. Cross-border demand, $D_i$, should increase with $P_i$ and decrease with distance $d_i$:
\begin{equation}
D_i = D(P_i,d_i),
\qquad
\frac{\partial D_i}{\partial P_i} > 0,
\qquad
\frac{\partial D_i}{\partial d_i} < 0.
\end{equation}
The intuition is simple. A larger nearby population creates a larger pool of potential cross-border
consumers, while greater distance raises the cost of reaching a legal dispensary.

Cross-border demand adds to the local demand already faced by dispensaries in county $i$. Let $R_i$ denote
total cannabis revenue and $N_i$ the number of active dispensaries. If dispensary entry adjusts freely, an
increase in demand may lead to more stores rather than much higher revenue at existing stores. This follows
the standard market-size and entry logic in \citet{BresnahanReiss1991}. Cannabis markets, however, often
face licensing, zoning, and other local restrictions that can limit entry \citep{HollenbeckUetake2021}. When
the number of dispensaries cannot adjust quickly, additional cross-border demand is more likely to raise
revenue per existing dispensary:
\begin{equation}
\frac{R_i}{N_i}
=
f\left(D_i(P_i,d_i),\,L_i,\,N_i\right),
\end{equation}
where $L_i$ represents local demand. Under constrained entry,
\begin{equation}
\frac{\partial (R_i/N_i)}{\partial P_i} > 0
\end{equation}
and revenue per dispensary should fall as distance from the nonlegal market increases.

This framework maps directly onto our three results. Result 1 tests whether revenue per dispensary is higher
closer to a nonlegal border, using both a binary and a continuous version of $d_i$, and whether the pattern
is absent in a sector without a comparable legal discontinuity. Result 2 tests the supply-constraint
condition directly: a state without a binding license cap should show the border effect in revenue per
capita rather than revenue per dispensary, since entry can adjust to absorb demand. We test this using
Montana in Section~\ref{sec:mechanism}. Result 3 tests whether the effect is stronger where $P_i$ is larger,
which distinguishes a genuine captive-demand channel from an unrelated correlate of border proximity.

\section{Data}
\label{sec:data}

Cannabis revenue comes from each state's own tax administrative data. For Colorado we use the Department of
Revenue's county-level marijuana sales series on the Colorado Information Marketplace
\citep{COMarijuanaSalesByCounty}, which reports monthly medical and recreational sales by county from 2014
onward; we aggregate to calendar years and keep only county-years with all twelve months reported. For
Washington we use the Liquor and Cannabis Board's fiscal-year sales and excise tax reports by county
\citep{WALCBSalesByCounty}, using the retailer-only sales figures in each report. Retail sales began January
1, 2014 in Colorado \citep{COMedLicensees} and July 8, 2014 in Washington \citep{WSLCB2014FirstLicenses}.
Dispensary counts are built from each state's own current licensee registry rather than a geocoded historical
listing: the Colorado Marijuana Enforcement Division's licensee lookup tool \citep{COMedLicensees} and the
Washington Liquor and Cannabis Board's own current retailer list, which reports each licensee's county
directly rather than requiring an address to be geocoded and matched to a county polygon
\citep{WALCBLicenseeList,WALCBCannabisRetailers2026}, applied as a current snapshot to each county's panel
years, consistent with how the per-store outcome is constructed throughout. Restricting the resulting
county-year panel to observations with at least one open dispensary leaves $n=787$ county-years: 42 of
Colorado's 64 counties and 38 of Washington's 39 counties, the remainder never having an operating dispensary
over the sample period. Population comes from the U.S. Census Bureau's Population Estimates Program, county
population totals, 2023 vintage \citep{CensusPopEst2023}; we use each county's July 1, 2023 estimate
throughout the panel rather than a year-specific series, consistent with treating population as a
slow-moving control rather than a time-varying outcome.

A county is classified as a border county if its Census Bureau internal point, the standard representative
coordinate the Bureau reports for each county rather than a population-weighted centroid, lies within 70
miles of the internal point of the nearest county in a state without an operating recreational retail
market. This classification is fixed for each county over the full sample period; it is built from the most
recent map of state legalization dates available to the project rather than updated year by year as
individual neighboring states legalize, so within-panel changes in border status do not occur by
construction. The continuous distance measure used in Result 1 is the natural log of this same distance in
miles; no specification we report uses inverse distance or distance bins. $P_i$, the potential nonlegal-side
market used in Result 3, is the combined population of every county in a nonlegal state within 70 miles of
county $i$, not only the nearest one; it uses the same fixed 2023 population vintage and the same
nonlegal-state map as the border classification and is therefore also fixed over time.

In Colorado, 6 of the 42 sample counties meet the border definition, bordering Utah, Wyoming, or Nebraska,
and the remaining 36 are interior. Because dispensaries did not open in every county at once, the number of
Colorado counties observed in the panel grows from 25 in 2015 to 41 from 2023 onward; border counties grow
correspondingly from 3 in 2015 to the full set of 6 from 2020 onward, averaging 5.3 border and 30.4 interior
counties per year. In Washington, 6 of the 38 sample counties are border counties, all near Idaho; this count
grows from 4 in 2015 and 2016 to 5 in 2017 and the full set of 6 from 2018 onward, averaging 5.6 border and
30.4 interior counties per year. Table~\ref{tab:sample_construction} summarizes these patterns by year. At
the observation level, the main Colorado-Washington regression sample of $n=787$ county-years splits into
119 border and 668 interior observations (58 and 334 in Colorado, 61 and 334 in Washington).

Montana serves as the contrast case in Result 2 because it imposes no statewide numerical cap on dispensary
licenses. Its adult-use market was established by House Bill 701, signed May 18, 2021 \citep{MTHB701}, with
retail sales beginning January 1, 2022 \citep{MTRevenueCannabis}; retail sales are permitted only in counties
where a majority of voters approved the underlying 2020 ballot initiative. We use the Department of Revenue's
monthly county-level cannabis sales reports \citep{MTRevenueCountySales} and its licensed-dispensary list
\citep{MTDispensaryLocations}. A county-year enters the sample only if all twelve months are disclosed;
Montana's confidentiality rule suppresses county-months with too few licensed operators to protect individual
businesses, and suppression is more common in smaller counties, so the number of counties with complete data
falls from 35 in 2022 and 2023 to 24 and 23 in 2024 and 2025. This is a reporting artifact of the disclosure
rule, not evidence of market contraction. Of the counties with complete data, 11 meet the 70-mile border
definition in 2022 and 2023 (near Idaho, North Dakota, or Wyoming), falling to 9 in 2024 and 8 in 2025 as
smaller border counties are suppressed, averaging 9.8 border and 19.5 interior counties per year over the
four-year sample; at the observation level, the Montana regression sample of $n=117$ county-years splits into
39 border and 78 interior observations.

The general-retail placebo used in Result 1 uses the Colorado Department of Revenue's county-by-industry
retail sales filings \citep{CDORRetailByCounty}. This source discloses industry-level detail only for the 11
Colorado counties with the largest filing volumes, a count we confirm directly from the underlying filings
rather than from the source's own description; the resulting sample of $n=110$ county-years splits into 30
border and 80 interior observations.

\begin{table}[htbp]
\centering
\caption{Sample Construction: Treated and Control Counties by Year}
\label{tab:sample_construction}
\small
\begin{tabular}{lcccl}
\toprule
State & Years & Border counties & Interior counties & Non-legal neighboring states \\
\midrule
Colorado & 2015 & 3 & 22 & UT, WY, NE \\
Colorado & 2016 & 4 & 25 & UT, WY, NE \\
Colorado & 2017--2019 & 5 & 27--30 & UT, WY, NE \\
Colorado & 2020--2025 & 6 & 31--35 & UT, WY, NE \\
Washington & 2015--2016 & 4 & 26--29 & ID \\
Washington & 2017 & 5 & 29 & ID \\
Washington & 2018--2025 & 6 & 31--32 & ID \\
Montana & 2022--2023 & 11 & 24 & ID, ND, WY \\
Montana & 2024 & 9 & 15 & ID, ND, WY \\
Montana & 2025 & 8 & 15 & ID, ND, WY \\
\bottomrule
\end{tabular}
\begin{minipage}{0.95\textwidth}
\vspace{0.3em}
\footnotesize
\textit{Notes:} Border and interior counts are the counties actually observed in each state's revenue panel
in that year, not the full statewide county count. Colorado's and Washington's border definitions are fixed
per county for the full sample; year-to-year changes reflect counties entering the panel as they obtain
their first dispensary, not changes in border status. Montana's decline in 2024--2025 reflects the state's
confidentiality-suppression rule rather than market contraction; see text.
\end{minipage}
\end{table}

\section{Empirical Strategy}

We estimate
\begin{equation}
\ln\left(\frac{R_{ist}}{N_{ist}}\right) = \alpha + \beta_1\,\text{Border}_i + \beta_2 \ln(P_i) + \gamma\ln(\text{Population}_{ist}) + \theta_t + \delta_s + \varepsilon_{ist},
\end{equation}
where $\text{Border}_i$ is a binary indicator for counties within 70 miles of a nonlegal border (or,
separately, $\ln(d_i)$ as a continuous measure, in Result 1), $P_i$ is the captive nonlegal population within
70 miles (Result 3), and $\theta_t$ and $\delta_s$ are year and state fixed effects. We report
heteroskedasticity-robust (HC1) and county-clustered standard errors throughout, and treat the clustered
estimates as our preferred inference given the panel structure. We estimate the same specification for
revenue per capita to test which margin responds (Results 1 and 2).

\section{Result 1: Does a Border Premium Exist?}
\label{sec:result1}

\begin{table}[htbp]
\centering
\caption{Result 1: The Border Premium in Revenue per Dispensary}
\label{tab:result1}
\small
\begin{tabular}{p{6.6cm}cccc}
\toprule
Specification & Coef. & SE (HC1) & SE (cl.) & $p$ (cl.) \\
\midrule
\multicolumn{5}{l}{\textit{Binary: border ($\le$70mi) vs.\ interior}} \\
No controls & $+0.382$ & $0.099$ & $0.219$ & $0.080$ \\
Full controls$^{a}$ & $+0.341$ & $0.075$ & $0.161$ & $0.033$ \\
\quad excl.\ Sedgwick County, CO & $+0.302$ & $0.068$ & $0.147$ & $0.040$ \\
Colorado only & $+0.552$ & -- & $0.199$ & $0.005$ \\
Washington only & $+0.233$ & -- & $0.243$ & $0.334$ \\
\midrule
\multicolumn{5}{l}{\textit{Continuous: log(distance to nonlegal border, miles)}} \\
Full controls$^{a}$ & $-0.215$ & $0.048$ & $0.095$ & $0.023$ \\
\quad excl.\ Sedgwick County, CO & $-0.168$ & $0.048$ & $0.092$ & $0.067$ \\
\bottomrule
\end{tabular}
\begin{minipage}{0.95\textwidth}
\vspace{0.3em}
\footnotesize
\textit{Notes:} Outcome is log(revenue per active dispensary). County-year panel, Colorado and Washington,
$n=787$ (80 counties: 42 Colorado, 38 Washington; 12 border counties: 6 Colorado, 6 Washington; 68 interior).
$^{a}$Full controls: log(population), year fixed effects, state fixed effects. Both the binary and
continuous measures are significant under county-clustered inference; a county closer to a nonlegal border
earns significantly more revenue per dispensary, and the relationship holds using a continuous distance
measure rather than only the 70-mile cutoff.
\end{minipage}
\end{table}


Table~\ref{tab:result1} reports the main result. Counties within 70 miles of a nonlegal border earn
significantly higher revenue per dispensary: 41 percent higher with the full control set, county-clustered
$p=0.033$. This is not driven by a single county; excluding Sedgwick County, Colorado, a small, heavily
touristed county on the Nebraska border, the estimate falls only to $+0.302$ and remains significant. It
holds within Colorado alone ($p=0.005$); the Washington-specific estimate is directionally consistent but
not itself significant under clustering ($p=0.334$) once Washington's dispensary counts are measured
correctly. The continuous distance measure tells the same story and is the more robust of the two: it is
significant under both HC1 and county-clustered inference ($p=0.023$), which is what a genuine
distance-decay relationship should look like rather than an artifact of the 70-mile cutoff.

\begin{table}[htbp]
\centering
\caption{Result 1 Robustness: Which Margin Responds, and a Non-Cannabis Placebo}
\label{tab:result1_placebo}
\small
\begin{tabular}{p{7.4cm}cc}
\toprule
Outcome & Coefficient (border effect) & $p$-value (cl.) \\
\midrule
log(revenue per active dispensary) & $+0.341$ & $0.033$ \\
log(revenue per capita) & $+0.046$ & $0.833$ \\
\midrule
\multicolumn{3}{l}{\textit{Placebo: general retail sales per capita, Colorado, 2016--2025}} \\
Border, full controls & $-0.006$ & $0.947$ \\
\bottomrule
\end{tabular}
\begin{minipage}{0.95\textwidth}
\vspace{0.3em}
\footnotesize
\textit{Notes:} The revenue-per-dispensary premium in Result 1 does not appear in total per-capita revenue,
consistent with the license-cap mechanism in Result 2. The general-retail placebo uses Colorado Department
of Revenue county-industry filings for the 11 counties with disclosed industry detail
\citep{CDORRetailByCounty}, $n=110$ (30 border, 80 interior); the precise null supports a cannabis-specific
rather than a generic border-commerce interpretation of Result 1.
\end{minipage}
\end{table}


Table~\ref{tab:result1_placebo} shows the premium is specific to revenue per dispensary. Total revenue per
capita shows no comparable border pattern, and general retail sales in Colorado, a sector with no legal
discontinuity at the same border, show a precise, well-powered null ($p=0.947$), arguing against a general
border-commerce explanation.

\subsection{A Direct Test: Colorado's Border with New Mexico}
\label{sec:co_nm}

The comparisons above are cross-sectional: border status never changes for a given county, so they cannot
rule out a stable, unobserved difference between border and interior counties that happens to look like a
demand effect. New Mexico's own retail market opened on April 1, 2022, partway through our panel, switching
six Colorado counties (Archuleta, Conejos, Costilla, La Plata, Las Animas, and Montezuma) from bordering a
nonlegal state to bordering a legal one. If the premium above reflects genuine captive cross-border demand,
it should fall in these counties once New Mexico's own market can absorb that demand, exactly the mechanism
\citet{HansenMillerWeber2020} document for Washington and Oregon, but on a border and event they do not
study.

\begin{table}[htbp]
\centering
\caption{Natural Experiment: Colorado's Border with New Mexico}
\label{tab:co_nm}
\small
\begin{tabular}{lcccc}
\toprule
Specification & Coef.\ (NM-border $\times$ post-2022) & SE (HC1) & SE (cl.) & $p$ (cl.) \\
\midrule
Difference-in-differences$^{a}$ & $-0.487$ & $0.106$ & $0.123$ & $<0.001$ \\
\bottomrule
\end{tabular}
\begin{minipage}{0.95\textwidth}
\vspace{0.3em}
\footnotesize
\textit{Notes:} Outcome is log(revenue per active dispensary). Colorado county-year panel, 2014--2025, $n=352$
(2022, the transition year in which New Mexico's market opened partway through, is dropped). Treatment is the
6 Colorado counties within 70 miles of New Mexico (Archuleta, Conejos, Costilla, La Plata, Las Animas,
Montezuma); controls include the 5 Colorado counties bordering a state that remained non-legal throughout the
sample (Utah, Wyoming, Oklahoma, Nebraska, Kansas) and the 31 Colorado interior counties, plus
log(population) and year fixed effects. $^{a}$Post $=1$ for 2023 onward, after New Mexico's April 1, 2022
retail opening. The implied change in revenue per dispensary in New Mexico-border counties after New Mexico's
own market opens is $-38.6$ percent.
\end{minipage}
\end{table}


\begin{figure}[htbp]
\centering
\includegraphics[width=0.75\textwidth]{figures/co_nm_eventstudy.png}
\caption{Event-study path for Colorado counties bordering New Mexico, coefficient relative to 2021, around
New Mexico's April 1, 2022 retail opening (event year 0). Shaded regions are 95\% confidence intervals using
county-clustered standard errors.}
\label{fig:co_nm}
\end{figure}

Table~\ref{tab:co_nm} reports the result. Revenue per dispensary in New Mexico-border counties falls 38.6
percent relative to Colorado's other border and interior counties after New Mexico opens its own market
(2023 onward, with 2022 dropped as a transition year; county-clustered $p<0.001$). Figure~\ref{fig:co_nm}
shows the same pattern as an event-study path against two separate comparison groups: relative to Colorado's
other nonlegal-state-border counties, which holds ``being a border county'' fixed and isolates only New
Mexico's change in legal status, the gap is not significantly different from zero in any pre-period year, and
it opens exactly at the 2022 switch and widens monotonically through 2025, the pattern a genuine
demand-switching mechanism predicts and a stable, unobserved difference between border and interior counties
could not produce on its own.

We attempted the mirror-image replication of \citet{HansenMillerWeber2020}'s own setting, Washington counties
bordering Oregon around Oregon's October 2015 retail opening, using our county-level annual data; that
replication does not cleanly confirm their result, and we report it in Appendix~\ref{app:wa_or} rather than
omit it. We view the Colorado-New Mexico result above, not the Washington-Oregon attempt, as this paper's
genuine causal evidence for the captive-demand mechanism.

\section{Result 2: Evidence of License-Cap Policy}
\label{sec:mechanism}

The pattern in Table~\ref{tab:result1_placebo} is exactly what the framework in Section~3 predicts if entry
is supply-constrained: additional demand cannot translate into more dispensaries, so it appears as higher
revenue per existing store instead. Colorado and Washington both constrain retail entry. Washington capped
the statewide number of licenses at 334 from the outset, allocating them by lottery in oversubscribed
jurisdictions \citep{WSLCB2014LotteryPR}; Colorado imposes no statewide cap but gives local governments full
authority under Amendment 64 to ban retail marijuana outright or to regulate its number, and most Colorado
municipalities have exercised this option to prohibit it entirely.

\begin{table}[htbp]
\centering
\caption{Result 2: Evidence of License-Cap Policy, Capped vs.\ Uncapped States}
\label{tab:result2}
\small
\begin{tabular}{lcc}
\toprule
Outcome & Coefficient (border effect) & $p$-value (county-cl.) \\
\midrule
\multicolumn{3}{l}{\textit{Colorado + Washington (license caps bind)}} \\
log(revenue per dispensary) & $+0.341$ & $0.033$ \\
log(revenue per capita) & $+0.046$ & $0.833$ \\
\midrule
\multicolumn{3}{l}{\textit{Montana (no statewide numerical cap)}} \\
log(revenue per capita), 2022--2025 panel & $+0.654$ & $0.036$ \\
log(revenue per dispensary), 2025 cross-section & $+0.012$ & $0.958^{a}$ \\
\bottomrule
\end{tabular}
\begin{minipage}{0.95\textwidth}
\vspace{0.3em}
\footnotesize
\textit{Notes:} $^{a}$HC1 standard error; too few Montana border counties ($n=23$) for county-clustered
inference. Colorado and Washington constrain entry (Washington through a statewide license cap allocated by
lottery; Colorado through local governments' authority under Amendment 64 to ban or limit retail marijuana).
Montana imposes no equivalent statewide cap. Where entry is capped, the border effect appears in revenue per
dispensary and not per capita; where entry is not capped, it appears in revenue per capita and not per
dispensary, consistent with a license-cap explanation for why the border premium in Result 1 shows up on the
intensive rather than the extensive margin.
\end{minipage}
\end{table}


Montana offers a within-sample test of this logic, since it is the one state in our data with no statewide
numerical cap on dispensary licenses. If the license-cap mechanism is right, Montana's border counties should
show the opposite pattern: captive demand absorbed through more stores (higher revenue per capita) rather
than through higher revenue at each existing store. Table~\ref{tab:result2} shows exactly this reversal:
Montana's border effect is significant in revenue per capita ($p=0.036$) and a precise null in revenue per
store ($p=0.958$). A single coefficient could be a coincidence; a predicted reversal that shows up in exactly
the state whose institutional structure differs in exactly the relevant way is harder to dismiss as one.

\section{Result 3: How Is Cross-Border Demand Captured?}
\label{sec:result3}

Results 1 and 2 establish that a border premium exists and responds to license-cap policy as predicted, but
neither speaks to whether the premium is genuinely about the size of the captive market on the nonlegal
side, rather than some other correlate of being close to a state line.

\begin{table}[htbp]
\centering
\caption{Result 3: Does the Premium Scale with the Nonlegal-Side Population?}
\label{tab:result3}
\small
\begin{tabular}{p{7.4cm}cccc}
\toprule
Specification & Coef. & SE (HC1) & SE (cl.) & $p$ (cl.) \\
\midrule
log($P_i$: captive nonlegal population, 70mi) + full controls & $+0.030$ & $0.007$ & $0.014$ & $0.036$ \\
\quad excl.\ Sedgwick County, CO & $+0.024$ & $0.007$ & $0.014$ & $0.078$ \\
\bottomrule
\end{tabular}
\begin{minipage}{0.95\textwidth}
\vspace{0.3em}
\footnotesize
\textit{Notes:} Outcome is log(revenue per active dispensary), $n=787$. $P_i$ sums the population of
\emph{all} nonlegal-state counties within 70 miles of county $i$, not only the nearest one, so it
distinguishes a county next to a small nonlegal county from a county next to a large nonlegal metropolitan
area, something the distance measures in Result 1 cannot do. Only 10 of the 80 Colorado and Washington
counties in this panel have any captive nonlegal population within 70 miles at all, so this result reweights
the same identifying set of border counties from Result 1 by market size rather than adding new geographic
variation; combined with Result 1's continuous-distance finding, it indicates the premium responds
continuously to both how close a county is to a nonlegal border and how large the population on the other
side is, not only to a binary border classification.
\end{minipage}
\end{table}


Table~\ref{tab:result3} tests this directly using $P_i$, the combined population of all nonlegal-state
counties within 70 miles of county $i$, not only the nearest one. Revenue per dispensary rises with $P_i$,
and this result survives county-clustered inference ($p=0.036$). We view this, together with the
continuous-distance result in Result 1, as the strongest evidence that Result 1's border premium reflects a
genuine captive cross-border demand channel: a county next to one small nonlegal county and a county next to
a large nonlegal metropolitan area are treated identically by a nearest-county distance measure, but only
the latter has a large population of potential cross-border shoppers, and only the latter shows the larger
premium Result 3 predicts.

\section{Robustness}
\label{sec:robustness}

Result 1 is robust to excluding the largest outlier county and to estimating Colorado and Washington
separately. The general-retail placebo addresses the concern that border counties might simply transact more
retail business of any kind, for reasons unrelated to cannabis policy; the precise null there is evidence
against that alternative explanation. Result 3's captive-population measure addresses a different concern,
that the 70-mile threshold is an arbitrary cutoff, by showing the relationship is graded in the size of the
nonlegal-side market rather than a discontinuous jump at 70 miles alone. We note directly that the
Washington-specific estimate in Result 1 is not significant under clustering once Washington's dispensary
counts are corrected from an earlier, incomplete registry; we report the corrected, weaker estimate rather
than the more favorable uncorrected one. The New Mexico natural experiment in Section~\ref{sec:co_nm}
addresses a different concern than these robustness checks: rather than showing the cross-sectional premium
survives alternative specifications, it shows the premium responds to an actual change in a neighboring
state's legal status, which a stable, unobserved difference between border and interior counties cannot
explain.

\section{Discussion}

These results speak to both state and local policy. For nonlegal states, the results imply a measurable
fiscal externality: some of the revenue and associated economic activity that would otherwise be available
to a nonlegal state if it legalized and taxed recreational cannabis itself instead accrues to legal-state
border counties. Our estimates do not measure the counterfactual tax revenue a nonlegal state would collect
if it legalized; they measure how existing legal-state dispensary revenue varies with proximity to, and the
population of, nearby nonlegal markets.

For legal states, the results imply that licensing policy in border counties is not purely a local question.
If border-county dispensaries serve a substantially larger effective market than local population alone
would suggest, a licensing cap set with only local population in mind may be too restrictive relative to
actual demand, generating exactly the revenue-per-dispensary premium documented here without necessarily
reflecting a well-functioning competitive market, a concern directly connected to the market-power findings
in \citet{HollenbeckUetake2021}.

Several limitations should be kept in view. Our revenue data do not identify the state of residence of
individual consumers, so most of our evidence on the demand channel is indirect, built from distance,
population-exposure, and placebo patterns rather than direct data on consumer origin; the New Mexico natural
experiment in Section~\ref{sec:co_nm} is the exception, since it identifies the demand channel from an actual
change in a neighboring state's legal status. Our core analysis for Results 1 and 3 uses only two states with
usable county-level revenue and dispensary data; Result 2's Montana contrast, while informative, rests on a
single additional state and a shorter panel. Apart from the New Mexico natural experiment, the paper's
identification relies primarily on cross-sectional and within-state variation in border proximity rather than
a policy-timing event, since a county's border status does not otherwise change over time; the results outside
Section~\ref{sec:co_nm} should therefore be read as strong, mechanism-consistent associations rather than as
fully causal estimates.

\section{Conclusion}

We ask three questions about cannabis retail markets near U.S.\ state borders. First, does a border premium
exist: yes, revenue per active dispensary is significantly higher in border counties than interior counties,
using both a binary and a continuous distance measure, and it does not appear in a non-cannabis placebo
sector; a natural experiment around New Mexico's 2022 legalization shows the premium falls 38.6 percent in
Colorado's New Mexico-border counties once New Mexico's own market can absorb the demand, direct causal
evidence for the mechanism rather than only a cross-sectional association. Second, is this evidence of
license-cap policy: yes, the premium appears on the per-dispensary margin
in Colorado and Washington, which restrict entry, and reverses to the per-capita margin in Montana, which
does not. Third, how is this demand captured: the premium scales with the size of the nonlegal-side
population within reach, not merely with proximity to it. Building on \citet{HansenMillerWeber2020} and
\citet{HaoCowan2020}, who each establish cross-border cannabis effects in a single setting, this paper shows
the pattern generalizes across multiple state borders, scales with the size of the market on the other side,
and is directly attributable to a specific, testable feature of state cannabis policy.

\appendix

\section{Robustness Check: Replicating Hansen, Miller, and Weber (2020) with County-Level Annual Data}
\label{app:wa_or}

As a robustness check on the New Mexico natural experiment in Section~\ref{sec:co_nm}, we attempted the
mirror-image test in \citet{HansenMillerWeber2020}'s own setting: Washington counties bordering Oregon around
Oregon's October 1, 2015 retail opening, using our own county-level annual revenue and dispensary
construction rather than their monthly transaction data. We compare Washington's 13 Oregon-border counties
(within 70 miles of Oregon) to Washington's 4 Idaho-border counties and its remaining interior counties,
using fiscal years 2015 through 2025 (FY2015 is the only fully pre-legalization year available).

\begin{table}[htbp]
\centering
\caption{Replication Attempt: Washington's Border with Oregon}
\label{tab:wa_or}
\small
\begin{tabular}{lcccc}
\toprule
Specification & Coef. & SE (HC1) & SE (cl.) & $p$ (cl.) \\
\midrule
Transition year only (FY2016 vs.\ FY2015) & $-0.545$ & $0.329$ & $0.336$ & $0.105$ \\
DiD, post $=$ FY2016 onward & $+0.136$ & $0.381$ & $0.432$ & $0.753$ \\
DiD, FY2016 dropped, post $=$ FY2017 onward & $+0.161$ & $0.380$ & $0.443$ & $0.716$ \\
\bottomrule
\end{tabular}
\begin{minipage}{0.95\textwidth}
\vspace{0.3em}
\footnotesize
\textit{Notes:} Outcome is log(revenue per active dispensary). Washington county-year panel, fiscal years
2015--2025, $n=395$ (2015 is the only fully pre-legalization year). Treatment is the 13 Washington counties
within 70 miles of Oregon; controls include the 4 Idaho-border counties, the remaining interior counties,
log(population), and fiscal-year fixed effects. The transition-year estimate is directionally consistent with
\citet{HansenMillerWeber2020} but only marginally significant; the longer-run DiD estimates are positive and
insignificant, the opposite of sustained demand suppression. See Appendix~\ref{app:wa_or} for discussion.
\end{minipage}
\end{table}


\begin{figure}[htbp]
\centering
\includegraphics[width=0.75\textwidth]{figures/wa_or_eventstudy_corrected.png}
\caption{Event-study path for Washington counties bordering Oregon, coefficient relative to fiscal year 2015,
around Oregon's October 1, 2015 retail opening (inside fiscal year 2016). Shaded region is the 95\%
confidence interval using county-clustered standard errors.}
\label{fig:wa_or}
\end{figure}

This replication does not cleanly confirm \citet{HansenMillerWeber2020}'s result. The immediate transition
year (fiscal year 2016, roughly three pre-legalization and nine post-legalization months) shows a
directionally consistent dip of similar sign to their finding, though only marginally significant under
clustering ($p=0.10$), but the coefficient reverses to significantly positive from fiscal year 2017 onward and
remains elevated through 2025 (Figure~\ref{fig:wa_or}), the opposite of the sustained suppression a genuine
demand-switching effect would predict. We attribute this to two features of our design rather than to the
underlying mechanism being wrong: our Oregon-border county group is dominated by large, fast-growing
metropolitan counties (Clark County, population 521,150, directly across the Columbia River from Portland;
Thurston, population 299,003; and Benton, population 215,219) whose revenue trends plausibly reflect general
population and urban growth rather than the 2015 switch specifically, and only one full pre-period fiscal
year is available, not enough to check for pre-existing divergence before treatment.
\citet{HansenMillerWeber2020} avoid both problems with store-level monthly transaction data and a
difference-in-discontinuities design; our coarser annual, county-level data are simply not suited to this
particular replication. We report this null result rather than omit it, and rely on the New Mexico natural
experiment in Section~\ref{sec:co_nm}, whose treated counties are small and rural rather than large and
rapidly urbanizing, for this paper's causal evidence on the demand-switching mechanism.

\section{Exploratory Analysis: Entry Density and Fiscal Spillover}
\label{app:welfare}

As a secondary analysis outside the paper's three main results, we tested whether the license-cap mechanism
in Result 2 also appears on the entry-density margin directly, that is, whether Colorado and Washington
border counties have systematically fewer dispensaries per capita than interior counties, which would be the
textbook signature of a binding entry restriction. Using the corrected Washington dispensary data described
in Section~\ref{sec:data}, this result does not hold: the Colorado-Washington border coefficient on
dispensaries per capita is small and statistically indistinguishable from zero ($b=-0.05$, $p=0.46$), a
materially different finding from an earlier version of this analysis that relied on an undercounted
Washington dispensary registry and had found a significant negative effect. We report this null directly
rather than omit it. Result 2's per-capita versus per-dispensary contrast between Colorado, Washington, and
Montana remains the paper's evidence for the license-cap mechanism; this entry-density test was intended as
additional, direct evidence on the same mechanism, and it does not corroborate it once Washington's
dispensary data are corrected.

Separately, scaling Result 1's estimated border coefficient ($b=0.341$, a 41 percent premium) by Colorado and
Washington border counties' actual 2025 revenue (\$351 million across 10 counties under the corrected border
definition) implies approximately \$101 million in revenue attributable to the border effect that year. We
present this as an accounting decomposition of an already-estimated coefficient, not as an independent
causal magnitude or a welfare calculation. Higher revenue per dispensary under a binding cap is
observationally consistent with either dispensaries serving a genuinely larger customer base or exploiting
the local market power that restricted entry creates, and \citet{HollenbeckUetake2021} document exactly this
kind of market power for Washington's licensing regime generally; our revenue data contain no prices or
quantities, so we cannot adjudicate between these explanations, and we do not characterize this figure as a
welfare effect.

\newpage
\bibliographystyle{apalike}
\bibliography{references}

\end{document}
